% ============================================================================
% instrument.tex -- test harness for the \package renderer
% ============================================================================
%
% \package typesets its output; a test can only assert on what reaches the log.
% Rather than dump boxes (unreadable, and sensitive to fonts and TeX Live
% version), we replace the three leaves of the renderer with log markers:
%
%   \texttt{NAME}      -> NAME=[...]
%   \href{URL}{ICON}   -> LINK=[...]
%   \faCode            -> (nothing)
%
% So `\package*{JAX}' logs
%
%   NAME=[JAX]
%   LINK=[https://pypi.org/project/jax]
%
% and an unregistered `\package*{nosuch}' logs a NAME line and *no* LINK line,
% which is precisely the fallback behaviour we want to pin down.  A wrong URL,
% a missing link, or a link that should not be there is then a one-line diff in
% the golden .tlg.
%
% Load this AFTER preamble-base-starkman (it redefines that package's leaves),
% and only in tests: it destroys the real rendering.
% ============================================================================

\renewcommand{\texttt}[1]{\typeout{NAME=[#1]}}
\renewcommand{\href}[2]{\typeout{LINK=[#1]}}
\renewcommand{\faCode}{}

% Show the resolved link for NAME without going through the renderer at all,
% i.e. assert directly against the registry rather than against \package.
% Logs `MISS=[NAME]' when NAME is not registered.
% NB: these must be defined here, in expl3 syntax, rather than written inline in
% a test: the body of \TEST is tokenized under normal catcodes, where `_' and
% `:' are not letters and \prop_get:NnNTF is not a control sequence.
\ExplSyntaxOn

\NewDocumentCommand \SHOWLINK { m }
  {
    \prop_get:NnNTF \g_starkman_packagelink_prop {#1} \l_tmpa_tl
      { \typeout { REGISTERED=[\l_tmpa_tl] } }
      { \typeout { MISS=[#1] } }
  }

% Dump the entire lookup table, keys and values.
\NewDocumentCommand \SHOWREGISTRY { }
  { \prop_show:N \g_starkman_packagelink_prop }

\ExplSyntaxOff
